\documentclass[12pt]{article}
\usepackage[margin=0.75in]{geometry}
\usepackage{fontspec}
\setmainfont{Latin Modern Roman}
\usepackage{microtype}
\usepackage{multicol}
\usepackage{fancyhdr}
\usepackage{titlesec}
\usepackage{enumitem}
\usepackage{xcolor}
\usepackage[hidelinks]{hyperref}
\setlength{\columnsep}{0.28in}
\setlength{\parindent}{1.1em}
\setlength{\parskip}{0.32em}
\setlength{\headheight}{15pt}
\titleformat{\section}{\large\bfseries\scshape}{}{0pt}{}
\titleformat{\subsection}{\normalsize\bfseries}{}{0pt}{}
\pagestyle{fancy}
\fancyhf{}
\fancyfoot[C]{\thepage}
\renewcommand{\headrulewidth}{0.4pt}
\newcommand{\focusbox}[1]{\par\vspace{0.4em}\noindent\begingroup\setlength{\fboxsep}{6pt}\colorbox{gray!12}{\begin{minipage}{\dimexpr\linewidth-2\fboxsep\relax}#1\end{minipage}}\endgroup\par\vspace{0.5em}}

\fancyhead[L]{Whately: Disputes}
\fancyhead[R]{Student Reading}
\begin{document}
\begin{center}
{\LARGE\bfseries Verbal and Real Disputes}\par\vspace{0.25em}
{\large\scshape Week 15: Are We Actually Disagreeing?}\par\vspace{0.5em}
{Adapted and modernized from Richard Whately, \textit{Elements of Logic}}\par\vspace{0.6em}
\rule{0.82\textwidth}{0.4pt}
\end{center}

\section*{Central Question}
How can we tell whether people disagree about reality or only about words?

\section*{Vocabulary Preview}
\begin{multicols}{2}
\begin{itemize}[leftmargin=*, itemsep=0.18em]
\item \textbf{Dispute}: a disagreement that needs to be tested for its real point.
\item \textbf{Verbal dispute}: a disagreement caused mainly by different uses of a word.
\item \textbf{Real dispute}: a disagreement that remains even after words are defined.
\item \textbf{Mixed dispute}: a disagreement involving both word-confusion and real disagreement.
\item \textbf{Definition}: the meaning a speaker gives to a key term.
\item \textbf{Contested term}: a word whose meaning is part of the disagreement.
\item \textbf{Shared claim}: something both sides would accept once terms are clarified.
\item \textbf{Remaining disagreement}: what the sides still dispute after clarification.
\item \textbf{Dispute map}: contested term, each side's meaning, shared claim, remaining disagreement, classification.
\item \textbf{Charitable clarification}: stating both sides fairly before judging them.
\end{itemize}
\end{multicols}

\section*{Introduction}
Many arguments are harder than they need to be because people do not first ask what kind of disagreement they are having. Two people may fight over a word while agreeing about the facts. Two others may define the word clearly and still disagree deeply. A third argument may contain both: confused words and a real dispute underneath.

Whately treats this distinction as a serious part of logic because clear thinking often begins by finding the exact question. If the problem is verbal, define the term and much of the quarrel may disappear. If the problem is real, definitions will not settle it; evidence and argument are still needed.

\focusbox{\textbf{Source note:} Adapted and modernized from Whately's treatment of verbal and real questions. The emphasis is not that all disputes are ``just semantics.'' The goal is to separate word-confusion from genuine disagreement so students can argue about the real point. Classroom examples are original.}

\begin{multicols}{2}
\section*{Reading}

\subsection*{1. Not Every Disagreement Is the Same}
When people disagree, they may be doing different things:
\begin{itemize}[leftmargin=*]
\item using the same word in different senses;
\item disagreeing about facts;
\item disagreeing about value or priority;
\item agreeing more than they realize;
\item disagreeing both about words and about reality.
\end{itemize}
A disciplined thinker does not rush to win the argument. First, identify what the argument is actually about.

\subsection*{2. Verbal Disputes}
A verbal dispute happens when the conflict is mainly about the meaning of a word. Suppose two students argue, ``This assignment is long.'' One means it has many pages. The other means it will take a long time. Once they clarify the word \textit{long}, they may discover that both are right in different senses: the assignment is only two pages but requires slow thinking.

The dispute was not empty. It showed a real ambiguity. But the solution is definition, not more shouting.

\subsection*{3. Real Disputes}
A real dispute remains after the key terms are clear. Suppose two students agree that by \textit{fair} they mean ``the rule gives each student the same chance to succeed,'' but one thinks a late-work policy does that and the other thinks it does not. Now the disagreement is not merely verbal. They need evidence, examples, principles, and judgment.

Definitions can prepare a real dispute, but they do not decide it by themselves.

\subsection*{4. Mixed Disputes}
Many serious arguments are mixed. A school may debate whether a punishment is \textit{just}. One side may mean ``allowed by the rules.'' Another may mean ``proportionate to the offense.'' Clarifying the word helps, but it may not end the disagreement. The class may still disagree about what the rule requires, what proportion is fair, or what facts matter.

Mixed disputes require two steps: clarify the term, then identify what remains.

\subsection*{5. The Danger of ``Just Semantics''}
Some people dismiss verbal disputes by saying, ``That is just semantics.'' That can be lazy. Words guide thought. If the same word carries two meanings, the argument can move without anyone noticing.

But the opposite mistake is also common: pretending a real disagreement will vanish if people define words. Sometimes people understand each other and still disagree. Logic should reduce confusion without pretending that every conflict is fake.

\subsection*{6. How to Map a Dispute}
Use this dispute map:
\begin{enumerate}[leftmargin=*,itemsep=0.12em]
\item State the disputed claim.
\item Identify the contested term.
\item State what Side A means by the term.
\item State what Side B means by the term.
\item Name any shared claim after clarification.
\item Name what remains disputed.
\item Classify the dispute as verbal, real, or mixed.
\end{enumerate}

\subsection*{7. Literature and Public Speech}
In \textit{Macbeth}, characters fight over words such as \textit{man}, \textit{safe}, and \textit{fate}. In \textit{Julius Caesar}, speakers fight over \textit{honor}, \textit{ambition}, \textit{liberty}, and \textit{tyranny}. These are not dictionary games. Characters use contested words to justify action, shame others, persuade crowds, and hide from the real point.

A disciplined reader asks: Is the conflict verbal, real, or mixed? If the word were clarified, what would still have to be proved?

\section*{Reading Focus}
\begin{enumerate}[leftmargin=*]
\item What is a verbal dispute?
\item What is a real dispute?
\item Why are many serious disputes mixed?
\item Why can ``just semantics'' be a lazy response?
\item What remains to be done after a disputed term is clarified?
\item Copy the seven steps of the dispute map.
\item Name one word from \textit{Macbeth} or \textit{Julius Caesar} that might create a mixed dispute.
\end{enumerate}

\section*{Handout Summary}
Write 5--7 sentences explaining how to distinguish verbal, real, and mixed disputes. Your summary must include the words \textit{contested term}, \textit{definition}, \textit{shared claim}, and \textit{remaining disagreement}.
\end{multicols}
\end{document}
